%----------------------------------------------------------------------------------------
% NOTATIE EN FIGUREN VAN DE CURSUS GROEPEN
% Staat in de preambule en niet in een hoofdstuk: ook de lijsten met oplossingen gebruiken dit.
%----------------------------------------------------------------------------------------
% Bewaarde grafiek, geplaatst door maak-pdf.py: bestand en pagina in figuren-cache/
\newcommand{\grafiekbestand}[2]{\includegraphics[page=#2,trim=3bp 3bp 3bp 3bp]{figuren-cache/#1}}

% Getallenverzamelingen en de gehele getallen modulo n
\newcommand{\N}{\mathbb{N}}
\newcommand{\Z}{\mathbb{Z}}
\newcommand{\Q}{\mathbb{Q}}
\newcommand{\R}{\mathbb{R}}
\newcommand{\C}{\mathbb{C}}
\newcommand{\Zn}[1]{\mathbb{Z}_{#1}}
% Optellen en vermenigvuldigen modulo n
\newcommand{\plusn}{\oplus}
\newcommand{\maaln}{\odot}
% Een permutatie in twee rijen: \perm{1 & 2 & 3 \\ 2 & 3 & 1}
\newcommand{\perm}[1]{\begin{pmatrix} #1 \end{pmatrix}}

% Zoals stappen, maar met een bredere kolom voor de berekening: voor lange formules
\newenvironment{stappenbreed}{%
  \par\renewcommand{\arraystretch}{1.25}%
  \noindent\begin{tabular}{@{}>{$\displaystyle}p{9.8cm}<{$}@{\hspace{4mm}}>{\RaggedRight\small\itshape\arraybackslash}p{4.9cm}@{}}%
}{\end{tabular}\par}

% Zoals stappen, met eigen kolombreedtes: \begin{stappenvrij}{6.2cm}{3.8cm}. Voor naast een figuur.
\newenvironment{stappenvrij}[2]{%
  \par\renewcommand{\arraystretch}{1.25}%
  \noindent\begin{tabular}{@{}>{$\displaystyle}p{#1}<{$}@{\hspace{3mm}}>{\RaggedRight\small\itshape\arraybackslash}p{#2}@{}}%
}{\end{tabular}\par}

% Tekst of berekening links, figuur rechts: \naast[breedte rechts]{links}{rechts}
% De breedte rechts komt uit de figuur zelf; het optionele argument wordt niet gebruikt.
\newsavebox{\naastbox}
\newcommand{\naast}[3][]{%
  \par\noindent\sbox{\naastbox}{#3}%
  \begin{minipage}[t]{\dimexpr\linewidth-\wd\naastbox-3mm\relax}\vspace{0pt}#2\end{minipage}\hfill
  \begin{minipage}[t]{\wd\naastbox}\vspace{0pt}\usebox{\naastbox}\end{minipage}\par\addvspace{1mm}}
% Een rij figuren, gecentreerd
\newcommand{\figuren}[1]{\par\addvspace{1mm}{\centering #1\par}\addvspace{1.5mm}}
% Opgave met een grafiek rechts: \begin{metgrafiek}[breedte]{tikz} opgave en deelvragen \end{metgrafiek}
% De solution komt na \end{metgrafiek}, zoals bij metprent.
\newsavebox{\grafiekbox}
\newenvironment{metgrafiek}[2][]{%
  \sbox{\grafiekbox}{#2}%
  \begin{minipage}[t]{\dimexpr\linewidth-\wd\grafiekbox-3mm\relax}%
}{%
  \end{minipage}\hfill
  \begin{minipage}[t]{\wd\grafiekbox}\vspace{-2mm}\usebox{\grafiekbox}\end{minipage}\par}

% BEWERKINGSTABEL, met een raster en alles gecentreerd. Alle vakjes staan in wiskundemodus.
% \begin{bewtabel}[breedte vakje]{aantal elementen}{bewerking}{kopregel}
%   a & a \ast a & a \ast b \\ \hline
%   ...
% \end{bewtabel}
% Rij a, kolom b: daar staat a * b. Een vakje kleuren: \cellcolor{fel!35}.
\newenvironment{bewtabel}[4][6.5mm]{%
  \renewcommand{\arraystretch}{1.3}%
  \begin{tabular}{|>{\columncolor{accent!14}$}wc{#1}<{$}|*{#2}{>{$}wc{#1}<{$}|}}\hline
  \rowcolor{accent!14} #3 & #4 \\ \hline
}{\end{tabular}}

% FIGUREN, te gebruiken binnen een tikzpicture
% Gelijkzijdige driehoek met straal 1 en de nummers van de hoekpunten:
% \driehoek[middelpunt]{boven}{linksonder}{rechtsonder}
\newcommand{\driehoek}[4][0,0]{%
  \begin{scope}[shift={(#1)}]
  \draw[accent,thick,fill=accent!8,line join=round] (90:1) -- (210:1) -- (330:1) -- cycle;
  \node[font=\small,inner sep=1pt] at (90:1.3) {$#2$};
  \node[font=\small,inner sep=1pt] at (210:1.3) {$#3$};
  \node[font=\small,inner sep=1pt] at (330:1.3) {$#4$};
  \end{scope}}
% De drie symmetrieassen van \driehoek: \driehoekassen[middelpunt]
\newcommand{\driehoekassen}[1][0,0]{%
  \begin{scope}[shift={(#1)}]
  \foreach \h in {90,210,330}{\draw[grijs,dashed,thin] (\h:1.15) -- (\h+180:.75);}
  \end{scope}}
% Vierkant met de nummers van de hoekpunten, tegen de wijzers van de klok in vanaf rechtsboven:
% \vierkant[middelpunt]{rechtsboven}{linksboven}{linksonder}{rechtsonder}
\newcommand{\vierkant}[5][0,0]{%
  \begin{scope}[shift={(#1)}]
  \draw[accent,thick,fill=accent!8,line join=round] (.75,.75) -- (-.75,.75) -- (-.75,-.75) -- (.75,-.75) -- cycle;
  \node[font=\small,inner sep=1pt] at (1,1) {$#2$};
  \node[font=\small,inner sep=1pt] at (-1,1) {$#3$};
  \node[font=\small,inner sep=1pt] at (-1,-1) {$#4$};
  \node[font=\small,inner sep=1pt] at (1,-1) {$#5$};
  \end{scope}}
% Rechthoek (geen vierkant), zoals \vierkant: \rechthoek[middelpunt]{rechtsboven}{linksboven}{linksonder}{rechtsonder}
\newcommand{\rechthoek}[5][0,0]{%
  \begin{scope}[shift={(#1)}]
  \draw[accent,thick,fill=accent!8,line join=round] (1.1,.6) -- (-1.1,.6) -- (-1.1,-.6) -- (1.1,-.6) -- cycle;
  \node[font=\small,inner sep=1pt] at (1.35,.82) {$#2$};
  \node[font=\small,inner sep=1pt] at (-1.35,.82) {$#3$};
  \node[font=\small,inner sep=1pt] at (-1.35,-.82) {$#4$};
  \node[font=\small,inner sep=1pt] at (1.35,-.82) {$#5$};
  \end{scope}}
% Klok voor het rekenen modulo n: \klok[straal]{n}, met 0 bovenaan en de getallen met de wijzers mee
\newcommand{\klok}[2][1]{%
  \draw[grijs,thick] (0,0) circle (#1);
  \foreach \k in {1,...,#2}{%
    \pgfmathsetmacro{\hoek}{90-360*(\k-1)/#2}%
    \pgfmathtruncatemacro{\getal}{\k-1}%
    \fill[accent] (\hoek:#1) circle (1.7pt);
    \node[font=\small,inner sep=1pt] at (\hoek:#1+.3) {$\getal$};}}
% Pijlendiagram van een permutatie: \pijlen{2,3,1} (het beeld van 1, 2, 3, ...): bovenaan x, onderaan het beeld
\newcommand{\pijlen}[1]{%
  \foreach \b [count=\a] in {#1}{%
    \node[font=\small,inner sep=1.5pt,fill=white] at (.85*\a,1) {$\a$};
    \node[font=\small,inner sep=1.5pt,fill=white] at (.85*\a,0) {$\a$};
    \draw[-{Stealth[length=1.8mm]},accent,thick,shorten >=6pt,shorten <=6pt] (.85*\a,1) -- (.85*\b,0);}}
% Een draaipijl om het middelpunt: \draaipijl[middelpunt]{straal}{label}
\newcommand{\draaipijl}[3][0,0]{%
  \draw[fel!85!black,thick,-{Stealth[length=2mm]}] ($(#1)+(30:#2)$) arc[start angle=30,end angle=150,radius=#2];
  \node[font=\small,inner sep=1pt] at ($(#1)+(90:#2+.22)$) {#3};}
